\documentclass[a4paper,10pt,twocolumn]{article}
\usepackage[latin1]{inputenc}
\usepackage[english]{babel}
\usepackage{amsmath}
\usepackage{amsfonts}
\usepackage{amssymb}
\usepackage{titling}
\usepackage{nomencl}
\usepackage{booktabs}
\usepackage{color}
\usepackage[font={footnotesize}]{caption}
\usepackage{graphicx}
\usepackage[style=ieee,backend=bibtex]{biblatex}
\usepackage{fancyhdr}
\usepackage{varioref}
\usepackage[activate={true,nocompatibility},final,tracking=true,kerning=true,spacing=true,factor=1100,stretch=10,shrink=10]{microtype}
% activate={true,nocompatibility} - activate protrusion and expansion
% final - enable microtype; use "draft" to disable
% tracking=true, kerning=true, spacing=true - activate these techniques
% factor=1100 - add 10% to the protrusion amount (default is 1000)
% stretch=10, shrink=10 - reduce stretchability/shrinkability (default is 20/20)

% Reduce tracking around small caps to 40%
\SetTracking{encoding={*}, shape=sc}{40}

% Path to images.
\graphicspath{{img/}}

% Setup nomenclature.
\makenomenclature
\newlength{\nomtitlesep}
\setlength{\nomtitlesep}{\nomitemsep}
\setlength{\nomitemsep}{-\parsep}
\renewcommand{\nomgroup}[1]{%
    \itemsep\nomtitlesep
    \ifthenelse{\equal{#1}{V}}{\item[\textbf{Variables}]}{}
    \itemsep\nomitemsep
}

% Setup bibiliography.
\addbibresource{bibliography}

% Document info.
\author{Z0966990}
\title{Statistics Assignment}
\date{\today}

% Header and footer.
\pagestyle{fancy}
\fancyhf{}
\lhead{\thetitle}
\rhead{\theauthor}
\cfoot{\thepage}
\renewcommand{\headrulewidth}{0pt}
\renewcommand{\footrulewidth}{0pt}

\begin{document}
    
% Title page.
\begin{titlepage}
    \centering
    \vspace*{\fill}
    \includegraphics[width=0.5\textwidth]{Durham}\\
    \vspace*{\fill}
    \LARGE\thetitle\\
    \large\theauthor\\
    \large L2 Engineering Mathematics\\
    \large\thedate\\
    \vspace*{\fill}
\end{titlepage}

% Nomenclature.
\nomenclature[0]{$n$}{Number of meter readings in sample.}
\nomenclature[1]{$n_{out}$}{Number of meter readings outside the 95\% 
confidence interval for the mean.}
\nomenclature[2]{$\bar{x}$}{Mean total energy consumption for sample.}
\nomenclature[3]{$\sigma$}{Unbiased estimate of population standard deviation.}

\nomenclature[V]{$E$}{Total energy consumption.}
\nomenclature[V]{$P$}{Rate of energy consumption.}
\nomenclature[V]{$t$}{Time.}
\printnomenclature

% Main matter.
\section{Introduction}

Statistical analysis was used to identify the demand patterns of energy 
consumption in a range of homes on a summer and winter's day, but also to 
analyse the effect of outliers on a distribution.

In 2009 the Irish Commission for Energy Regulation (CER) undertook a trial 
where smart meters were installed in 6256 homes \cite{martinsmart}. Those 
meters recorded the rate of electrical energy consumption every half hour for a 
year.

This report analysed the total energy consumption for two days in the 
trial---the first day which occurred in the summer and the 200th day which 
occurred in the winter. In order to get an energy consumption in kWh, the  
following simple numerical integration was applied:

\begin{equation}
    E = \sum{P \cdot \Delta t}
    \label{eq:Integrate}
\end{equation}

Where $P$ was the the half hourly power demand readings in kW. Constant $\Delta 
t$ was 0.5~h, the time interval between readings. The variable analysed in this 
report ($E$) was an approximation of the energy consumed that day in kWh, 
henceforth denoted $x$.

It was noted a large number of meters recorded no energy consumption. This 
could have been due to the consumers being away from the premises or the meters 
being faulty. These meters were excluded from statistical analysis to avoid 
skewing the mean of the distribution.

\section{Demand Patterns}

Figure~\vref{fig:Data} shows the distribution of the energy consumption 
recorded by 6122 smart meters on day 1 and 6207 smarts meters day 200. The 
histograms do not include the meters which recorded no energy consumption that 
day.
    
\begin{figure}[h]
    \centering
    \def\svgwidth{0.5\textwidth}
    \input{img/Consumption.pdf_tex}
    \caption{Frequency histogram of total energy consumption on each day.}
    \label{fig:Data}
\end{figure}

From initial inspection, it was clear for both distributions the data had a 
strong negative skew with the majority of meters reporting a total consumption 
under 50~kWh, but with outliers as large as 500~kWh. This made it difficult to 
discern details in the dataset when only 50 bins were used.

Nonetheless, the mean total consumption still yielded statistically significant 
results. Table~\vref{table:Mean} records the upper and lower bounds for the 
mean at two different confidence levels.

\begin{table*}[h]
    \centering
    \caption{Confidence intervals for mean total energy consumption on each 
        day.\vspace{-\parsep}}
    \begin{tabular}{llll}
        \toprule
        \textbf{Day} & $\bar{x}$, \tiny{kWh} & 90\% Interval, \tiny{kWh} & 95\% 
        Interval, \tiny{kWh} \\
        \midrule
        \textbf{Summer} & 15.632 & [15.062, 16.203] & [14.952, 16.313] \\
        \textbf{Winter} & 19.788 & [19.127, 20.448] & [19.001, 20.575] \\
        \bottomrule
    \end{tabular}
    \label{table:Mean}
\end{table*}

The 90\% and 95\% confidence intervals for the population mean total energy 
consumption on each day did not overlap. The upper bound of the 95\% interval 
on the summer's day at 16.313~kWh was less than the lower bound of the 95\% 
interval on the winter's day at 19.001~kWh. This suggests, at the 5\% level of 
significance, the mean total energy consumption on day 200 in the winter for 
all Irish households was greater than on day 1 in the summer. The 90\% 
confidence intervals reiterate the same conclusion but with less significance.

\section{Outliers}

In order to eliminate the outliers in the dataset and achieve a more 
symmetrical distribution, the data for day 1 was revised by removing meters 
which recorded an energy consumption greater than the upper limit of the 95\% 
confidence interval for the mean. Figure~\vref{fig:NoOutliers} shows the 
distribution of the revised data.

\begin{figure}[h]
    \centering
    \def\svgwidth{0.5\textwidth}
    \input{img/Day1Revised.pdf_tex}
    \caption{Frequency histograms of original and revised total energy 
    consumption on the summer's day.}
    \label{fig:NoOutliers}
\end{figure}

From inspection, the distribution of the revised data appeared to follow a more 
normal distribution than the original data for energy consumption on the 
summer's day. This more normal distribution results in more accurate 
predictions for the 95\% confidence interval for the revised mean shown in 
Table~\vref{table:Outliers} alongside the original data. This is because the 
confidence intervals assume an underlying normal distribution.

\begin{table}[h]
    \centering
    \caption{Number of values outside the 95\% interval for mean total energy 
        consumption on the summer's day.\vspace{-\parsep}}
    \begin{tabular}{llll}
        \toprule
        \textbf{Data} & $\bar{x}$, \tiny{kWh} & 95\% Interval, \tiny{kWh} & 
        $n_{out}$ \\
        \midrule
        \textbf{Original} & 15.632 & [14.952, 16.313] & 3816 \\
        \textbf{Revised}  & 8.3173 & [8.2011, 8.4335] & 4595 \\
        \bottomrule
    \end{tabular}
    \label{table:Outliers}
\end{table}

Furthermore, the mean of the revised data lies closer to the modal and median 
values for the distribution. This suggests the mean of the revised data and its 
associated confidence intervals are a better measure of central tendency than 
the mean of the original data.

The number of data-points outside the 95\% confidence interval for the mean 
would be expected to decrease when the number of data-points was reduced in 
revising the original energy consumption data. On the contrary, 
Table~\vref{table:Outliers} indicates the number of points outside this 
interval increased from 3816/6122 to 4595/4660.

The width of a confidence interval is indicative of the precision in the 
estimate of the mean, given by $2\sigma/\sqrt{n}$. In this case, the reduction 
in variance due to removing outliers had a greater effect on increasing 
precision than reducing the sample size had on decreasing precision. In 
revising the energy consumption data, the width of the 95\% confidence interval 
dropped from 1.541~kWh to 0.2324~kWh.

It was the increased precision in estimating the mean that resulted in more 
values lying outside the 95\% confidence interval.

% References.
\printbibliography

\clearpage

\end{document}